% ===========================================================================
% OWNER: M2 (Abdul Raffay)
% GRADED BY: Form 3 #7
% DRAW FROM: runner/README.md, runner/*.py, runner/harness/FtHarness.java, docs/04-Implementation/{sandbox-runner,diagnosis-runs,diagnose-cli}.md, docs/evidence-m2.md, ADR-003/004/005/007/008
% Rules: mid-eval-report/CLAUDE.md. Do not edit a file you do not own.
% ===========================================================================
\subsection{Ordered Runner, Polluter Search, Minimisation and CLI (Member 2)}\label{ch6:sec:m2}
Member 2's module is the \path{runner/} package: it runs tests in a chosen order, finds the polluter, shrinks the failing order, repeats it, and is the command that joins the three packages. It is 1{,}119 lines of Python and Java, with 1{,}488 lines of tests (\texttt{wc -l}, 11 October 2026).

\subsubsection{Files and main functions}
{\small
\begin{longtable}{>{\raggedright\arraybackslash}p{0.26\textwidth}>{\raggedright\arraybackslash}p{0.07\textwidth}>{\raggedright\arraybackslash}p{0.55\textwidth}}
\caption{Files of the runner module}\label{ch6:tab:m2-files}\\
\toprule
File & Lines & Job\\
\midrule
\endfirsthead
\toprule
File & Lines & Job\\
\midrule
\endhead
\path{order_runner.py} & 220 & \texttt{OrderRunner.run\_ordered}: one fresh JVM per order; \texttt{maven\_test\_classpath}; failure-signature normalisation\\
\path{harness/FtHarness.java} & 122 & Runs the listed tests in order with \texttt{JUnitCore}; lists test methods for discovery\\
\path{discovery.py} & 31 & \texttt{discover\_order}: test classes by Surefire's default include patterns, methods from JUnit\\
\path{search.py} & 55 & \texttt{reproduce} (first real failure becomes the reference); \texttt{find\_polluter} (one-by-one)\\
\path{orders.py} & 69 & \texttt{distinct\_shuffles}, \texttt{search\_orders}: class-first shuffled orders (ADR-008)\\
\path{minimise.py} & 55 & \texttt{ddmin}: 1-minimal polluter set (ADR-007)\\
\path{verify.py} & 33 & \texttt{repeat}: $n$ runs, matching and any-signature failure counts\\
\path{integrity.py} & 40 & \texttt{snapshot}, \texttt{compare}: SHA-256 of the project's files before and after\\
\path{recording.py} & 54 & \texttt{RecordingRunner}: writes every JVM run to the execution record\\
\path{diagnose.py} & 208 & \texttt{diagnose}, \texttt{run\_steps}: the whole diagnosis, returns raw counts\\
\path{cli.py}, \path{__main__.py} & 232 & \texttt{py -m runner diagnose}: options, evidence per polluter, report, summary, exit codes\\
\bottomrule
\end{longtable}}

\subsubsection{Ordered runner}
\texttt{OrderRunner} compiles \path{FtHarness.java} once with \texttt{javac -source 8 -target 8}. Each call to \texttt{run\_ordered(order)} then starts \texttt{java FtHarness <result-file>} in the analysed project's folder, with the project's own test classpath, sends the test list on standard input, and waits at most 120\,s. The harness runs each test with JUnit's \texttt{JUnitCore.run(Request.method(...))} and writes one line per test to the result file, so output printed by a test cannot corrupt the results. Python then turns each failure into a \texttt{FailureSignature}:
\begin{itemize}
  \item \textbf{exception type:} the class of the thrown exception (not its cause);
  \item \textbf{stack trace:} the frames from the throw point down to the first frame of JUnit's runner, Java reflection or the harness. Those lower frames differ between JDK 8 and JDK 21 and say nothing about the failure;
  \item \textbf{message:} the first line with numbers, hexadecimal ids, timestamps and paths masked. It is stored but not compared, because the same bug can print different values on each run.
\end{itemize}
The contract requires every test in the order to appear in the result, so the runner never drops one (Table~\ref{ch6:tab:m2-always}).

{\small
\begin{longtable}{>{\raggedright\arraybackslash}p{0.42\textwidth}>{\raggedright\arraybackslash}p{0.46\textwidth}}
\caption{How the runner reports every test (\texttt{runner/README.md})}\label{ch6:tab:m2-always}\\
\toprule
Situation & Reported as\\
\midrule
\endfirsthead
\toprule
Situation & Reported as\\
\midrule
\endhead
JVM runs longer than 120\,s & failure, \texttt{flaketrace.Timeout}\\
JVM exits before reporting a test (crash, \texttt{System.exit}) & failure, \texttt{flaketrace.JvmCrash}, exit code in the message\\
Class or method not found & failure with JUnit's own exception\\
\texttt{@Ignore} or a failed \texttt{Assume} & failure, \texttt{flaketrace.NotExecuted}: a skip never counts as a pass\\
The same test twice in one order & \texttt{ValueError} before any JVM starts\\
\bottomrule
\end{longtable}}

The \texttt{flaketrace.*} types mark failures the tool produced, not the test. Such a failure is counted as ``failed for any reason'' but never becomes the reference signature and never counts as a reproduction.

\subsubsection{Diagnosis flow}
\texttt{diagnose()} hashes the project, builds the classpath, discovers the tests, wraps the runner in a \texttt{RecordingRunner}, and calls \texttt{run\_steps} (Figure~\ref{ch5:fig:search-activity}). It returns a \texttt{DiagnosisRuns} object with raw counts and one of four statuses (Table~\ref{ch6:tab:m2-status}); it never decides the outcome, which is Member 3's \texttt{decide()}.

{\small
\begin{longtable}{>{\raggedright\arraybackslash}p{0.25\textwidth}>{\raggedright\arraybackslash}p{0.40\textwidth}>{\raggedright\arraybackslash}p{0.23\textwidth}}
\caption{Statuses returned by \texttt{diagnose()}}\label{ch6:tab:m2-status}\\
\toprule
Status & Meaning & Sequence repeated $n$ times\\
\midrule
\endfirsthead
\toprule
Status & Meaning & Sequence repeated $n$ times\\
\midrule
\endhead
\texttt{POLLUTER\_FOUND} & One earlier test, or a 1-minimal set found by \texttt{ddmin}, makes the victim fail with the reference signature & polluters + victim\\
\texttt{VICTIM\_FAILS\_ALONE} & The victim fails with nothing before it; no test is blamed & victim alone\\
\texttt{NO\_SINGLE\_POLLUTER} & A failure was seen but did not come back; no test is blamed & the failing order\\
\texttt{NOT\_REPRODUCED} & No real failure in $n$ runs of the order, in the shuffled orders, or in $n$ runs alone & the starting order\\
\bottomrule
\end{longtable}}

\subsubsection{Minimisation with ddmin}\label{ch6:sec:ddmin}
When no single earlier test makes the victim fail, \texttt{run\_steps} first runs the full order once more. Only if it fails again with the reference signature does \texttt{ddmin} (\path{runner/minimise.py}) shrink the tests before the victim; otherwise a one-off flaky failure would make every earlier test look guilty. Listing~\ref{ch6:lst:ddmin} is the core loop.

\begin{lstlisting}[language=Python,basicstyle=\footnotesize\ttfamily,columns=fullflexible,keepspaces=true,caption={Core of \texttt{ddmin} (\texttt{runner/minimise.py})},label=ch6:lst:ddmin]
current, k = list(prefix), 2
while len(current) > 1:
    chunks = _split(current, k)
    smaller = next((chunk for chunk in chunks if fails(chunk)), None)
    if smaller is not None:
        current, k = smaller, 2
        continue
    complements = ([t for t in current if t not in chunk] for chunk in chunks)
    smaller = next((rest for rest in complements if fails(rest)), None)
    if smaller is not None:
        current, k = smaller, max(k - 1, 2)
        continue
    if k >= len(current):
        break
    k = min(2 * k, len(current))
return current, runs
\end{lstlisting}

\texttt{fails(subset)} runs \texttt{subset + [victim]} in its original relative order in one fresh JVM and returns true only if the victim fails with the \emph{reference} signature; a crash or a different exception never removes a test. A cache keyed by the set of tests means no subset runs twice. The loop ends when \texttt{k} equals the set size, so its last round has tried removing each remaining test on its own: the result is \emph{1-minimal} (removing any one test makes the failure disappear), not necessarily the smallest set that exists.

On fixture F3 (run on 11 October 2026, Windows 11, JDK 21.0.9), the command printed \texttt{minimised: 14 earlier tests -> 2 polluters in 10 runs}, and the two polluters were the ground truth's \texttt{setFlagA} and \texttt{setFlagB}. The execution record of that run shows 65 JVM runs in total: 1 to reproduce, 20 alone, 14 for the one-by-one search, 10 for minimisation (including the one re-run of the full order) and 20 to verify.

\subsubsection{Verification and order search}
\texttt{repeat} (\path{runner/verify.py}) runs the final sequence $n$ times, each in a fresh JVM, and returns two counts: runs where the victim failed with the reference signature, and runs where it failed for any reason. The default is $n = 20$, agreed by all three members: 20 matching failures out of 20 give a Wilson lower bound of 0.839 and 19 out of 20 give 0.764, both above the 0.70 that Member 3's \texttt{decide()} requires for \texttt{VERIFIED} (\path{docs/contracts/interfaces.md}, question I3; ADR-004). The victim-alone check uses the same function on \texttt{[victim]}.

When the starting order never fails for real, \path{runner/orders.py} tries up to 31 distinct shuffled orders, one run each (ADR-008). Classes are shuffled first and then the methods inside each class, because JUnit never interleaves the methods of two classes; every seed is written to the execution record, so any shuffle can be replayed. In the worst case a diagnosis that never reproduces costs $20 + 31 + 20 = 71$ JVM runs. On the fixture, F4's discovered order never fails, and the failure was reproduced in the second shuffled order (seed 1); N3 never failed in 20 runs of its order, 31 shuffled orders or 20 runs alone, and its report gives the bound ``failure rate below 0.161 (95\% Wilson), not proof of reliability'' (\path{docs/evidence-m2.md}, 11 October 2026).

\subsubsection{Command line}
\texttt{py -m runner diagnose} (\path{runner/cli.py}, ADR-005) takes the options in Table~\ref{ch6:tab:m2-options} and maps every result to an exit code (Table~\ref{ch6:tab:m2-exit}). A known error prints one line, \texttt{error: ...}, with no stack trace; an unexpected bug still prints its trace, on purpose.

{\small
\begin{longtable}{>{\raggedright\arraybackslash}p{0.15\textwidth}>{\raggedright\arraybackslash}p{0.17\textwidth}>{\raggedright\arraybackslash}p{0.56\textwidth}}
\caption{Options of \texttt{py -m runner diagnose}}\label{ch6:tab:m2-options}\\
\toprule
Option & Default & Meaning\\
\midrule
\endfirsthead
\toprule
Option & Default & Meaning\\
\midrule
\endhead
\texttt{-{}-project} & required & Maven project folder, must contain \texttt{pom.xml}\\
\texttt{-{}-victim} & required & failing test as \texttt{Class\#method}\\
\texttt{-{}-n} & 20 & repeat count\\
\texttt{-{}-records} & \texttt{flaketrace-\allowbreak records} & folder for the record and report; refused if inside the project\\
\texttt{-{}-order} & none & file with the real failing order, one \texttt{Class\#method} per line\\
\texttt{-{}-shuffles} & 31 & shuffled orders to try if the order never fails; 0 turns it off\\
\texttt{-{}-seed} & 0 & first shuffle seed\\
\bottomrule
\end{longtable}}

{\small
\begin{longtable}{>{\raggedright\arraybackslash}p{0.06\textwidth}>{\raggedright\arraybackslash}p{0.82\textwidth}}
\caption{Exit codes of \texttt{py -m runner diagnose}}\label{ch6:tab:m2-exit}\\
\toprule
Code & Meaning\\
\midrule
\endfirsthead
\toprule
Code & Meaning\\
\midrule
\endhead
0 & Report written, whatever the outcome, including \texttt{UNRESOLVED}\\
2 & Input is wrong: victim not \texttt{Class\#method}, no \texttt{pom.xml}, $n < 1$, records folder inside the project, bad \texttt{-{}-order} file, victim not among the project's tests\\
1 & A tool failed: Maven, \texttt{java}, \texttt{javac}, a timeout during discovery, or the extractor\\
3 & The diagnosis ran but no report can be built: \texttt{NO\_SINGLE\_POLLUTER}\\
\bottomrule
\end{longtable}}

Listing~\ref{ch6:lst:m2-f1} is the summary printed for fixture F1 on 11 October 2026 (the two lines with the report and record paths are left out).
\begin{lstlisting}[basicstyle=\footnotesize\ttfamily,columns=fullflexible,keepspaces=true,breaklines=true,breakatwhitespace=true,caption={Summary for F1 (exit code 0)},label=ch6:lst:m2-f1]
VERIFIED  odfixture.ConfigVictimTest#expectsDefaultMode
  polluter:   odfixture.ConfigPolluterTest#pollute
  resource:   static-field odfixture.Config mode (write odfixture.ConfigPolluterTest#pollute@1 -> read odfixture.ConfigVictimTest#expectsDefaultMode@1)
  reproduced: 20/20 (lower bound 0.839)   alone: 0/20
\end{lstlisting}

\subsubsection{Testing}
The module has 144 tests in nine files under \path{runner/tests/}. Most of the logic (search, minimisation, order search, recording) is tested against Member 3's \texttt{FakeOrderRunner}, which returns scripted outcomes without a JVM and exists only under \path{eval/tests/}; the order runner, \texttt{diagnose()} and the command are also tested with real JVM runs on the fixture. Without \texttt{java} or Maven the real tests are skipped, except in CI, where \texttt{FLAKETRACE\_REQUIRE\_JVM} turns a skip into a failure. All 144 passed in CI on Linux with JDK 8 in 105.2\,s (pull request \#49) and locally on Windows with JDK 21 in 344.8\,s (Table~\ref{ch6:tab:setup-results}).

Most tests were written before the code they check and were seen failing first; key tests were also checked by breaking the code on purpose and watching them fail (\path{docs/evidence-m2.md} records both for each task). Two examples: changing the alone check from \texttt{alone\_successes >= 1} to \texttt{>= 100} made \texttt{test\_victim\_failing\_alone\_stops\_before\_the\_search} fail; and a test that reads a 20{,}001-line \texttt{-{}-order} file within 5\,s failed at 33.4\,s, which exposed a quadratic duplicate check, after which the same file took 0.094\,s.

\subsubsection{Limitations and a failure case}
\begin{itemize}
  \item \textbf{No sandbox.} The analysed project's tests run as ordinary processes with the user's permissions; there is no container yet. FlakeTrace should only be run on projects the user trusts.
  \item \textbf{JUnit 3 and 4 only.} JUnit 5 is not supported. Each method is its own JUnit request, so \texttt{@BeforeClass} and \texttt{@AfterClass} run once per method, not once per class as under Maven Surefire.
  \item \textbf{Signatures within one JDK.} Frames from the JDK itself above the cut keep their line numbers, which differ between JDK versions, so such signatures are only compared within one diagnosis. Only the top-level exception is compared, not its cause.
  \item \textbf{Discovery order.} Tests are discovered in alphabetical order; Surefire's default runs them in file-system order. For a failure seen in CI, the real order should be passed with \texttt{-{}-order}.
  \item \textbf{No report for \texttt{NO\_SINGLE\_POLLUTER}.} The report schema has no reason for it, so the command exits with code 3 and only the execution record remains.
  \item \textbf{31 shuffles is an assumption.} The figure comes from a study of Python projects~\cite{gruber2021flaky}.
\end{itemize}

\textbf{Failure case: N2 on one laptop.} N2 is meant to fail at random. On 9 October 2026 its real-JVM test failed: \texttt{diagnose()} returned \texttt{NOT\_REPRODUCED}, because N2 failed in none of 40 runs of its order and none of 40 runs alone, although a design experiment earlier that day had seen it fail 18 times in 60 runs alone. A probe showed that \texttt{System.nanoTime()} on that laptop only returned multiples of 100, so N2's random condition could never come true there; why the earlier run differed was not found. The diagnosis was right and the test's assumption was wrong: the test now checks that no polluter is ever blamed and accepts \texttt{NOT\_REPRODUCED} when N2 never fails (\path{docs/evidence-m2.md}, ADR-004).
